\documentclass[aspectratio=169, 11pt]{beamer}

% --- Packages & Setup ---
\usepackage[utf8]{inputenc}
\usepackage[T1]{fontenc}
\usepackage{helvet}
\renewcommand{\familydefault}{\sfdefault} % Times New Roman font for text and math
\usepackage{xcolor}
\usepackage{booktabs}
\usepackage{graphicx}
\usepackage{tikz}
\usetikzlibrary{calc}

\graphicspath{{../}{./}}

% --- Branding Colors ---
\definecolor{amritaMaroon}{HTML}{A4123F}

% --- Theme Setup ---
\usetheme{Madrid}
\usecolortheme[named=amritaMaroon]{structure}
\setbeamertemplate{navigation symbols}{}
\setbeamertemplate{footline}[frame number]

\begin{document}

% =========================================================================
% TITLE SLIDE (Research-Based Project)
% =========================================================================
\begin{frame}
    % Corner Box Indicator
    \begin{tikzpicture}[remember picture,overlay]
        \node[fill=amritaMaroon, text=white, font=\bfseries\footnotesize, anchor=north east, rounded corners=2pt, inner sep=5pt]
        at ($(current page.north east) + (-0.2cm,-0.2cm)$)
        {Type 1: Research Project};
    \end{tikzpicture}

    \centering
    \includegraphics[width=2.8cm]{logo.png}
    \vspace{0.1cm}

    {\large \textbf{\textcolor{amritaMaroon}{[RESEARCH PROJECT TITLE]}}} \\
    \vspace{0.05cm}
    {\footnotesize \textbf{23CSE498 -- Project Phase II (External Review)} \quad | \quad \textbf{Team Number: [GXXX]}} \\
    \vspace{0.15cm}

    % Student Table
    \begin{table}
        \centering
        \footnotesize
        \renewcommand{\arraystretch}{1.05}
        \begin{tabular}{|c|c|l|}
            \hline
            \textbf{S.No} & \textbf{Register Number} & \textbf{Student Name} \\
            \hline
            1 & [CB.SC.U4CSE23XXX] & [Student Name 1] \\
            \hline
            2 & [CB.SC.U4CSE23XXX] & [Student Name 2] \\
            \hline
            3 & [CB.SC.U4CSE23XXX] & [Student Name 3] \\
            \hline
            4 & [CB.SC.U4CSE23XXX] & [Student Name 4] \\
            \hline
        \end{tabular}
    \end{table}
    \vspace{0.1cm}

    {\footnotesize Guide: \textbf{[Guide Name]}, [Designation], Department of CSE} \\
    \vspace{0.1cm}
    {\scriptsize \today}
\end{frame}

% =========================================================================
% SECTION 1: PROBLEM DEFINITION, RESEARCH MOTIVATION & SOLUTION DESIGN
% =========================================================================
% =========================================================================
% SLIDE 2: GUIDE APPROVAL EMAIL CONFIRMATION (Mandatory)
% =========================================================================
\begin{frame}{Guide Review \& Submission Approval}
    \framesubtitle{Confirmation of Project Review and Approval from Guide}
    \centering
    \fbox{\begin{minipage}[c][5.4cm]{0.94\textwidth}
        \centering
        \textbf{[Insert Screenshot of Guide Approval Email Here]} \\[0.3cm]
        {\footnotesize [Mandatory: Attach a clear, readable screenshot of the email received from your project guide approving this presentation and final review submission]}
    \end{minipage}}
\end{frame}

\section{Problem \& Solution Design}

% --- SLIDE 1: PROBLEM DEFINITION & RESEARCH MOTIVATION ---
\begin{frame}{1. Problem Definition \& Research Motivation}
    \framesubtitle{Problem Statement, Identified Literature Gap \& Objectives}
    \begin{block}{Problem Statement}
        \begin{itemize}
            \item \textbf{Core Research Problem:} {[State the specific research problem and scope clearly]}
            \item \textbf{Domain Context:} {[Briefly state why this problem is important to solve]}
        \end{itemize}
    \end{block}

    \begin{block}{Existing Limitations \& Research Gap}
        \begin{itemize}
            \item \textbf{Summary of Existing Approaches:} {[Summarize how prior literature approaches this problem]}
            \item \textbf{Identified Research Gap:} {[Clearly state the specific limitation or gap addressed by this work]}
        \end{itemize}
    \end{block}

    \begin{block}{Research Objectives}
        \begin{itemize}
            \item \textbf{Objective 1:} {[State measurable research objective 1]}
            \item \textbf{Objective 2:} {[State measurable research objective 2]}
            \item \textbf{Objective 3:} {[State measurable research objective 3]}
        \end{itemize}
    \end{block}
\end{frame}

% --- SLIDE 2: PROPOSED METHODOLOGY & OVERALL WORKFLOW (FULL-SLIDE DIAGRAM) ---
% Note to students: Duplicate this frame if multiple workflow/architecture diagrams are required.
\begin{frame}{2. Proposed Methodology \& Overall Workflow}
    \framesubtitle{End-to-End Research Methodology / System Workflow Architecture}
    \centering
    \fbox{\begin{minipage}[c][5.2cm]{0.94\textwidth}
        \centering
        \textbf{[Insert Full-Width Proposed Methodology / Research Workflow Diagram]} \\[0.3cm]
        {\footnotesize [Placeholder: Include end-to-end stages from Input Data $\rightarrow$ Processing / Core Methodology $\rightarrow$ Output \& Evaluation]}
    \end{minipage}}
    \vspace{0.2cm}

    \begin{block}{Workflow Summary}
        \footnotesize {[Provide a concise 1--2 line walkthrough of the end-to-end methodology depicted above.]}
    \end{block}
\end{frame}

% --- SLIDE 3: METHODOLOGICAL COMPONENTS & FORMULATION ---
% Note to students: Duplicate this frame if additional slides are needed to explain sub-modules or formulations.
\begin{frame}{3. Methodological Components \& Formulation}
    \framesubtitle{Sub-Module Breakdown, Governing Mathematical Model \& Datasets}
    \begin{columns}[T]
        \column{0.50\textwidth}
        \begin{block}{Methodological Breakdown}
            \footnotesize
            \begin{itemize}
                \item \textbf{[Stage / Component 1]:} {[Role and technical function]}
                \item \textbf{[Stage / Component 2]:} {[Role and technical function]}
                \item \textbf{[Stage / Component 3]:} {[Role and technical function]}
            \end{itemize}
        \end{block}

        \begin{block}{Mathematical / Algorithmic Formulation}
            \footnotesize
            \begin{itemize}
                \item \textbf{Governing Formulation:} {[Insert key equation, objective function, or formal definition]}
                \item \textbf{Key Variables:} {[Define parameters and constraints]}
            \end{itemize}
        \end{block}

        \column{0.46\textwidth}
        \begin{block}{Dataset \& Input Specifications}
            \footnotesize
            \begin{itemize}
                \item \textbf{Source \& Size:} {[Dataset identifier, sample count, dimensions]}
                \item \textbf{Data Partitioning:} {[Splits / validation protocol used]}
                \item \textbf{Preprocessing:} {[Steps applied to raw inputs]}
            \end{itemize}
        \end{block}

        \begin{block}{Evaluation Metrics}
            \footnotesize
            \begin{itemize}
                \item \textbf{[Metric 1]:} {[Purpose / definition]}
                \item \textbf{[Metric 2]:} {[Purpose / definition]}
            \end{itemize}
        \end{block}
    \end{columns}
\end{frame}

% =========================================================================
% SECTION 2: IMPLEMENTATION COMPLETION (100%)
% =========================================================================
\section{Implementation Details}

% --- SLIDE 4: IMPLEMENTATION COMPLETION (100%) - OVERVIEW ---
\begin{frame}{4. Implementation Details -- Overview}
    \framesubtitle{Module-Wise Implementation \& Pipeline Integration Status}
    \begin{table}
        \centering
        \footnotesize
        \renewcommand{\arraystretch}{1.25}
        \begin{tabular}{p{3.6cm} p{6.2cm} c}
            \toprule
            \textbf{Module / Stage} & \textbf{Implemented Functionality} & \textbf{Status} \\
            \midrule
            {[Module 1: Data / Input Pipeline]} & {[Summary of completed data ingestion and preprocessing]} & 100\% (Integrated) \\
            {[Module 2: Baseline / Reference]} & {[Summary of implemented baseline or comparison setup]} & 100\% (Integrated) \\
            {[Module 3: Core Proposed Method]} & {[Summary of implemented core algorithm / model]} & 100\% (Integrated) \\
            {[Module 4: Evaluation \& Prototype]} & {[Summary of end-to-end execution and evaluation pipeline]} & 100\% (Functional) \\
            \bottomrule
        \end{tabular}
    \end{table}

    \begin{block}{Implementation Environment \& Repository}
        \footnotesize
        \begin{itemize}
            \item \textbf{Software / Compute Environment:} {[List programming environment, libraries, and compute resources used]}
            \item \textbf{Source Repository \& Execution Readiness:} {[Insert repository URL / commit reference and execution status]}
        \end{itemize}
    \end{block}
\end{frame}

% --- SLIDE 5: CORE MODULE IMPLEMENTATION DETAILS ---
% Note to students: Duplicate this frame if needed to explain additional implemented modules.
\begin{frame}{5. Core Module Implementation Details}
    \framesubtitle{Internal Module Mechanics, Parameter Configurations \& Interfaces}
    \begin{columns}[T]
        \column{0.50\textwidth}
        \begin{block}{[Core Module A Name]}
            \footnotesize
            \begin{itemize}
                \item \textbf{Inputs \& Outputs:} {[Specify module input and output formats]}
                \item \textbf{Implementation Details:} {[Explain how the module logic is realized in code]}
                \item \textbf{Key Parameters:} {[Specify configured parameters]}
            \end{itemize}
        \end{block}

        \begin{block}{[Core Module B Name]}
            \footnotesize
            \begin{itemize}
                \item \textbf{Inputs \& Outputs:} {[Specify module input and output formats]}
                \item \textbf{Integration Interface:} {[Explain how this module connects to the pipeline]}
            \end{itemize}
        \end{block}

        \column{0.46\textwidth}
        \centering
        \vspace{0.2cm}
        \fbox{\begin{minipage}[c][5.0cm]{0.94\textwidth}
            \centering
            \textbf{[Insert Module Interaction Diagram / Pipeline Execution Snapshot]} \\[0.2cm]
            {\scriptsize [Placeholder: Show how implemented modules exchange data or execute end-to-end]}
        \end{minipage}}
    \end{columns}
\end{frame}

% =========================================================================
% SECTION 3: PRELIMINARY VERIFICATION & DEMONSTRATION
% =========================================================================
\section{Verification \& Demonstration}

% --- SLIDE 6: PRELIMINARY VERIFICATION & QUANTITATIVE RESULTS ---
\begin{frame}{6. Preliminary Verification \& Quantitative Results}
    \framesubtitle{Experimental Benchmarking \& Comparative Evaluation}
    \begin{table}
        \centering
        \footnotesize
        \renewcommand{\arraystretch}{1.25}
        \begin{tabular}{p{4.5cm} c c c}
            \toprule
            \textbf{Method / Configuration} & \textbf{[Metric 1]} & \textbf{[Metric 2]} & \textbf{[Metric 3]} \\
            \midrule
            {[Baseline / Reference Method 1]} & {[Value]} & {[Value]} & {[Value]} \\
            {[Baseline / Reference Method 2]} & {[Value]} & {[Value]} & {[Value]} \\
            {[Ablation / Variant Configuration]} & {[Value]} & {[Value]} & {[Value]} \\
            \textbf{[Proposed Methodology]} & \textbf{[Value]} & \textbf{[Value]} & \textbf{[Value]} \\
            \bottomrule
        \end{tabular}
    \end{table}

    \begin{block}{Experimental Verification Summary}
        \footnotesize
        \begin{itemize}
            \item \textbf{Quantitative Findings:} {[Summarize key improvements or trends observed in the table above]}
            \item \textbf{Verification of Expected Behavior:} {[Explain how the results confirm the research hypothesis]}
        \end{itemize}
    \end{block}
\end{frame}

% --- SLIDE 7: EXPERIMENTAL VISUALIZATIONS & PROTOTYPE DEMONSTRATION ---
% Note to students: Use a single full-width box or two side-by-side boxes as appropriate for your plots/screenshots.
\begin{frame}{7. Experimental Visualizations \& Prototype Demonstration}
    \framesubtitle{Result Visualizations, Convergence / Performance Curves \& Live Demo}
    \begin{columns}[T]
        \column{0.48\textwidth}
        \centering
        \fbox{\begin{minipage}[c][4.8cm]{0.94\textwidth}
            \centering
            \textbf{[Insert Experimental Result Plot / Performance Graph]} \\[0.2cm]
            {\scriptsize [Placeholder: Quantitative chart, convergence curve, or comparative plot]}
        \end{minipage}}

        \column{0.48\textwidth}
        \centering
        \fbox{\begin{minipage}[c][4.8cm]{0.94\textwidth}
            \centering
            \textbf{[Insert Working Prototype / Qualitative Output Snapshot]} \\[0.2cm]
            {\scriptsize [Placeholder: Execution output, visualization, or demonstration screenshot]}
        \end{minipage}}
    \end{columns}
    \vspace{0.2cm}

    \begin{block}{Demonstration Overview}
        \footnotesize {[Briefly describe the live execution scenario or prototype walkthrough demonstrated to the panel.]}
    \end{block}
\end{frame}

% =========================================================================
% SECTION 4: TECHNICAL KNOWLEDGE & VIVA VOCE
% =========================================================================
\section{Technical Knowledge}

% --- SLIDE 8: TECHNICAL JUSTIFICATION, DESIGN TRADE-OFFS & INSIGHTS ---
\begin{frame}{8. Technical Justification, Design Trade-offs \& Insights}
    \framesubtitle{Methodological Rationale, Trade-offs \& Boundary Conditions}
    \begin{columns}[T]
        \column{0.48\textwidth}
        \begin{block}{1. Methodological Justification}
            \footnotesize
            \begin{itemize}
                \item {[Explain why the chosen theoretical/algorithmic approach is suited to the problem structure]}
                \item {[Justify key assumptions made in the formulation]}
            \end{itemize}
        \end{block}

        \begin{block}{2. Design \& Parameter Trade-offs}
            \footnotesize
            \begin{itemize}
                \item {[Discuss trade-offs between accuracy, computational cost, or convergence]}
                \item {[Justify selection of critical parameters]}
            \end{itemize}
        \end{block}

        \column{0.48\textwidth}
        \begin{block}{3. Boundary Conditions \& Edge Cases}
            \footnotesize
            \begin{itemize}
                \item {[Explain how noisy inputs, edge cases, or constraints are handled]}
                \item {[Summarize sensitivity or robustness observations]}
            \end{itemize}
        \end{block}

        \begin{block}{4. Limitations \& Future Scope}
            \footnotesize
            \begin{itemize}
                \item {[State current methodological or experimental limitations]}
                \item {[Outline concrete future extensions]}
            \end{itemize}
        \end{block}
    \end{columns}
\end{frame}

% =========================================================================
% REFERENCES & THANK YOU
% =========================================================================
\section{References}
\nocite{*}
\begin{frame}[allowframebreaks]{References}
    \bibliographystyle{ieeetr}
    \bibliography{ref}
\end{frame}

\begin{frame}
    \centering
    {\Huge \textbf{\textcolor{amritaMaroon}{Thank You}}} \\[0.4cm]
    {\large \textbf{[Questions \& Viva Voce]}}
\end{frame}

\end{document}
